% Generated draft: requires mathematical and source review.
\documentclass[11pt]{article}
\usepackage[margin=25mm]{geometry}
\usepackage{fontspec}
\setmainfont{Latin Modern Roman}
\usepackage{amsmath,amssymb,mathtools,unicode-math}
\setmathfont{Latin Modern Math}
\usepackage{enumitem,needspace,xurl,hyperref}
\setlength{\parindent}{0pt}
\setlength{\parskip}{5pt}
\setlength{\emergencystretch}{3em}
\newcommand{\N}{\mathbb{N}}
\newcommand{\Z}{\mathbb{Z}}
\newcommand{\Q}{\mathbb{Q}}
\newcommand{\R}{\mathbb{R}}
\newcommand{\C}{\mathbb{C}}
\newcommand{\F}{\mathbb{F}}
\newcommand{\A}{\mathbb{A}}
\newcommand{\PP}{\mathbb{P}}
\newcommand{\E}{\mathbb{E}}
\newcommand{\eps}{\varepsilon}
\newcommand{\dd}{\,\mathrm{d}}
\begin{document}
\section*{UnsolvedMath ID: 600015}
\textbf{Source problem code:} catalogue entry
\subsection{Definitions and mathematical statement}
Fix a combinatorial type $T$ of a dissection of a square into $n\geq1$ labelled triangles. A realization consists of a square partitioned into triangles with pairwise disjoint interiors and with the prescribed combinatorial data; let $a_i>0$ be the Euclidean area of triangle $i$. Write $\mathcal A_T$ for the resulting set of ordered area vectors $(a_1,\ldots,a_n)$. The precise incidence data and allowed degenerations constituting $T$ are not specified in the supplied record.

The stated premise is that there is a nonzero polynomial relation, depending only on $T$, satisfied by every vector in $\mathcal A_T$. The problem asks what can be said about this relation and how its least degree can be determined from $T$. For example, once a coefficient field and realization convention have been chosen, polynomial relations can be expressed as
\[
 I_T=\{f\in\Q[x_1,\ldots,x_n]: f(a_1,\ldots,a_n)=0
       \text{ for all }(a_1,\ldots,a_n)\in\mathcal A_T\}.
\]
A possible homogeneous version asks for
\[
 \min\{\deg f:0\ne f\in I_T,\ f\text{ homogeneous of positive degree}\}.
\]
This displays the distinction that must be resolved; it does not silently identify the source's ``least degree'' with this version.

In particular, if the square has fixed area $A$, every realization satisfies $\sum_i a_i=A$, an affine relation of degree one. The source's discussion of a distinguished irreducible homogeneous relation concerns a more specific area-polynomial formulation, whose normalization and treatment of this area-sum relation are not fully given. Allowing the square to scale, using normalized areas, and fixing a square lead to different ambient formulations. The intended relation ideal and degree notion must therefore be established before a unique algorithmic target is determined.
\subsection{Short English statement}
Determine the universal algebraic constraints on triangle areas for a prescribed square-dissection combinatorics, and compute the smallest degree of the intended relation from that combinatorics.

The registry labels the problem solved and its truncated research summary describes a distinguished irreducible homogeneous integer polynomial $p_T$, unique up to scalar, called the Monsky polynomial. It reports deformation invariance, a normalization with leading coefficients $\pm1$ in the parallelogram setting, and
\[
 p_T\equiv(x_1+\cdots+x_n)^d\pmod2,\qquad d=\deg p_T,
\]
as explaining Monsky's odd/even equidissection phenomenon. It also reports a principal relation ideal, a combinatorial degree lower-bound algorithm with exact worked examples, finiteness and explicit low-width relations denoted $\mathcal E_l$ for $l\leq4$, and direct enumeration for modest triangulation sizes. These are attributed claims of the supplied summary, not verified theorems asserted here: the numbered papers, width notation, normalization conventions, and complete final qualification are absent. The summary ends mid-sentence. The original 2015 source presents an open question, and the supplied status metadata explicitly requires review.
\subsection{Sources}
\begin{itemize}
\item UnsolvedMath permanent ID 600015, source code catalogue entry, ``Baker's Dozen --- Relations among triangle areas.'' Registry curation and metadata: Ulam AI and UnsolvedMath Contributors, CC BY 4.0; updated 2026-08-06.
\item Tabachnikov, \emph{A Baker's Dozen of Problems} (2015), Problem 7, supplied publisher HTML locator \url{https://amj.math.stonybrook.edu/html-articles/Files-2015-2024/14-01/}, accessed 2026-07-29 through an AMR-indexed list. Status evidence and rights review are marked NEEDS\_REVIEW in the record.
\item The research summary refers to ``paper 1'' through ``paper 5'' for its Monsky-polynomial and algorithmic claims, but supplies no authors, titles, or locators for those numbered references. No bibliography is invented for them.
\end{itemize}
\end{document}
